% Compile with pdfLaTeX. All article text and references are in this file.
% Independent article layout; not an official CSSS template.
\documentclass[11pt,letterpaper]{article}
\usepackage[T1]{fontenc}
\usepackage[utf8]{inputenc}
\usepackage{lmodern}
\usepackage[margin=0.9in,headheight=14pt,headsep=18pt]{geometry}
\usepackage{microtype}
\usepackage{parskip}
\usepackage{xcolor}
\usepackage{titlesec}
\usepackage{fancyhdr}
\usepackage{enumitem}
\usepackage{needspace}
\usepackage{xurl}
\usepackage{hyperref}
\definecolor{ink}{HTML}{193B4A}
\hypersetup{colorlinks=true,urlcolor=ink,linkcolor=ink,
 pdftitle={Who Pays for the Next Space Station?},
 pdfauthor={Adhvaith Ananth and Renish Joshi},pdfsubject={Space policy and economics}}
\urlstyle{same}
\setcounter{secnumdepth}{0}
\setlength{\parskip}{6pt}
\setlength{\emergencystretch}{2em}
\linespread{1.04}
\widowpenalty=10000
\clubpenalty=10000
\raggedbottom
\titleformat{\section}{\sffamily\bfseries\fontsize{26}{29}\selectfont\color{ink}}{}{0pt}{}
\titleformat{\subsection}{\sffamily\fontsize{16}{20}\selectfont\color{ink}}{}{0pt}{}
\titleformat{\subsubsection}{\sffamily\bfseries\large\color{ink}}{}{0pt}{}
\titlespacing*{\section}{0pt}{0pt}{7pt}
\titlespacing*{\subsection}{0pt}{0pt}{12pt}
\titlespacing*{\subsubsection}{0pt}{13pt}{5pt}
\setlist[enumerate]{leftmargin=1.5em,itemsep=6pt,parsep=0pt}
\providecommand{\tightlist}{\setlength{\itemsep}{0pt}\setlength{\parskip}{0pt}}
\pagestyle{fancy}
\fancyhf{}
\fancyhead[L]{\scriptsize\sffamily SPACE STATIONS AND PUBLIC SPENDING}
\fancyhead[R]{\scriptsize\sffamily Adhvaith Ananth and Renish Joshi}
\fancyfoot[C]{\small\thepage}
\renewcommand{\headrulewidth}{0.3pt}
\begin{document}
\hypertarget{who-pays-for-the-next-space-station}{%
\section{Who Pays for the Next Space
Station?}\label{who-pays-for-the-next-space-station}}

\hypertarget{government-customers-and-the-economics-of-commercial-orbit}{%
\subsection{Government customers and the economics of commercial
orbit}\label{government-customers-and-the-economics-of-commercial-orbit}}

\textbf{Adhvaith Ananth and Renish Joshi}\\
Policy and economics \textbar{} 6 October 2026

A privately owned space station still needs someone to pay for its use.
The difficult question is how much of that bill governments should
carry, and what they should receive in return.

\hypertarget{a-customer-as-well-as-a-sponsor}{%
\subsubsection{A customer as well as a
sponsor}\label{a-customer-as-well-as-a-sponsor}}

A space station needs more than a successful launch. Its operator must
keep equipment working, support the people aboard and arrange transport
for supplies and experiments. Those commitments continue between
customer visits. Before financing a station, an investor therefore needs
a plausible account of who will buy its services and how often.

NASA wants to buy services from commercial stations after the
International Space Station (ISS). Its stated ambition is to become one
customer among many, supporting research and exploration while a wider
market develops.{[}1{]} This raises a public spending question: how can
NASA help establish that market without making promises that taxpayers
must support indefinitely?

We argue that government support should be tied to identifiable public
needs, with separate assessments of technical progress and customer
demand. A station should not have to eliminate government revenue to
prove its worth. But private ownership alone should not be treated as
evidence that a broad commercial market exists. This article offers
policy analysis based on published documents, not a forecast of station
revenues or a financial assessment of individual companies.

\hypertarget{what-nasa-has-announced}{%
\subsubsection{What NASA has announced}\label{what-nasa-has-announced}}

NASA's plans have evolved. In June 2026, the Government Accountability
Office (GAO) reported uncertainty over the transition from the ISS and
warned that replacement stations might not be ready in time. It
recommended assessing the likelihood and duration of a gap and
documenting how NASA would decide whether to retire or extend the
ISS.{[}2{]}

A subsequent NASA announcement on 6 July described a draft procurement
approach using competitively awarded contracts. NASA said it intended to
select two or more contractors for early development, followed by
competition for later work, certification and services from one or more
providers.{[}3{]} That announcement describes an intended approach, not
proof of completed awards, operational stations or demonstrated demand.

The distinction matters when discussing public spending. Supporting a
design, certifying a station and purchasing its services are different
decisions. A successful design review may justify further development.
It does not establish how many customers will eventually pay to use the
finished station.

\hypertarget{follow-the-money-behind-the-customer}{%
\subsubsection{Follow the money behind the
customer}\label{follow-the-money-behind-the-customer}}

Consider a hypothetical university experiment booked through a private
station operator. The transaction looks commercial: a customer pays a
company for a service. Yet the experiment might be financed entirely by
a federal research grant. The station has gained a customer, but its
revenue still depends on public spending.

That arrangement can be worthwhile. Public research funding is meant to
support work whose benefits may extend beyond the organisation
conducting it. The accounting should simply make the source of demand
clear. Ten institutions spending government grants do not provide the
same independence from public budgets as ten businesses spending their
own revenue.

We suggest that programme assessments distinguish NASA purchases, other
government-funded activity and privately financed purchases.
International public customers could be identified separately where
disclosure is possible. Detailed customer contracts need not be
published; aggregate reporting could show the overall pattern while
protecting commercially sensitive information.

This would make claims about market development easier to evaluate. A
growing number of users might indicate broader scientific access,
greater private demand, or both. Each outcome has value, but they answer
different questions about the station's future financing.

\hypertarget{a-reliable-first-customer}{%
\subsubsection{A reliable first
customer}\label{a-reliable-first-customer}}

Government purchasing can make an uncertain project easier to finance.
If NASA identifies a service it expects to need and offers credible
terms for buying it, a provider can plan around something more concrete
than general enthusiasm for commercial spaceflight.

There is also a limit. Buying capacity simply to keep a company busy can
conceal weak demand. We would distinguish purchasing a service, paying
to keep essential capacity available, and supporting experimental
development. These activities should have explicit purposes and be
evaluated against those purposes. They should not all be presented as
routine purchases in a mature market.

For example, paying to preserve access to an orbital laboratory could be
defensible even during periods of low use. The justification would be
the value of maintaining that capability. Officials would still need to
explain how much availability they were buying, why it was necessary and
what alternatives were considered.

NASA's July announcement reports industry confidence in attracting
investment and developing additional markets.{[}3{]} That confidence is
relevant, but future purchasing decisions should also consider evidence
such as completed paid missions, repeat bookings and the conditions
attached to customer commitments. An expression of interest is a weaker
signal than a customer returning at a price that helps cover the
service's costs.

\hypertarget{count-the-whole-service}{%
\subsubsection{Count the whole service}\label{count-the-whole-service}}

A quoted price for laboratory space does not necessarily represent the
full cost of an orbital experiment. In a hypothetical purchase, the
customer may also need payload preparation, transport, crew assistance,
communications, data processing and sample return. Comparing providers
requires clarity about which services are included.

For government buyers, the same principle applies to comparisons with
the ISS. A claim of savings should identify the capability being
purchased and account for relevant transition and oversight costs. A
smaller research programme could cost less while delivering less. That
may be an acceptable budget choice, but it is different from obtaining
the same programme more efficiently.

Engineering choices influence these comparisons. An experiment designed
for remote operation may need less astronaut attention, while one
requiring frequent manual intervention may depend heavily on crew
availability. We recommend that buyers specify the service they need
closely enough to compare offers, while allowing providers room to
propose different technical solutions.

Standardised descriptions of power, data access and payload support
could also help researchers understand an offer. Standards would need to
be developed with users and operators; imposing one design too early
could exclude useful alternatives. Their purpose should be to reduce
avoidable integration work, not make every station identical.

\hypertarget{competition-after-the-first-award}{%
\subsubsection{Competition after the first
award}\label{competition-after-the-first-award}}

Funding several early designs can preserve choices while engineers learn
which approaches work. Maintaining several operating stations is a
different financial commitment. If the available demand is spread too
thinly, a policy intended to encourage competition could leave every
provider dependent on additional support.

The opposite problem deserves attention too. If NASA has only one
practical supplier, moving an experiment elsewhere may be difficult. A
nominal ability to compete contracts would offer limited protection if
no alternative could provide the necessary service within a useful time.

We therefore favour evaluating competition at each stage. Early
development should test credible alternatives. Later purchases should
consider price, reliability, switching costs and the value of backup
capacity. Supporting an additional provider might be justified for
resilience, but that expense should be identified rather than assumed to
pay for itself through competition.

Where feasible, portable experiment equipment and usable data exports
could reduce dependence on one operator. They would not make switching
effortless. Different station environments, schedules and safety reviews
would still matter. The aim is to avoid unnecessary barriers without
promising that orbital services can be exchanged as easily as ordinary
office space.

\hypertarget{research-deserves-its-own-justification}{%
\subsubsection{Research deserves its own
justification}\label{research-deserves-its-own-justification}}

An emphasis on private customers could disadvantage research with
substantial public value but limited immediate revenue. NASA identifies
crew training, research and hardware testing as reasons for continued
access to low Earth orbit.{[}1{]} Those purposes should be assessed
directly, rather than defended only through predictions of future
commercial products.

A public research purchase should explain why the orbital environment is
needed, how the work will be selected and what result will count as
useful. Some projects may produce openly available knowledge. Others may
appropriately involve protected commercial information. Access
conditions and expectations for results should be settled before the
work is funded.

Small research teams also need a practical route to participation. We
propose clear service descriptions, transparent application requirements
and support for preparing experiments. A station's availability means
little to a university team that cannot determine what a mission would
cost or how to obtain a usable result.

These measures would not guarantee demand. They would help reveal
whether a lack of bookings reflects limited interest in the research,
high total costs or avoidable difficulty in arranging a mission. Each
problem calls for a different response.

\hypertarget{plan-for-disappointing-demand}{%
\subsubsection{Plan for disappointing
demand}\label{plan-for-disappointing-demand}}

A credible purchasing strategy should remain workable if private demand
grows more slowly than expected. Before committing to a large volume of
services, officials should compare scenarios with strong private
participation, prolonged dependence on government customers and delays
in station availability.

Under weaker demand, possible responses could include reducing the
amount of capacity purchased, concentrating support on fewer services or
continuing a clearly justified public research programme. None is
automatically the right answer. Their costs and consequences should be
examined before a provider's financial difficulties make the decision
urgent.

Technical milestones and commercial milestones should remain separate. A
safe station could have too few customers; an attractive business plan
could depend on hardware that is not ready. Combining the two into a
single optimistic progress measure would make it harder to identify the
actual problem.

We also recommend periodic review of the public purpose behind
continuing support. Ending development funding immediately after a
missed target could waste a promising investment. Continuing it without
a defined reason could be equally wasteful. Reviews should explain what
has changed, what further support would purchase and when another
decision will be made.

\hypertarget{a-fair-test-of-commercial-success}{%
\subsubsection{A fair test of commercial
success}\label{a-fair-test-of-commercial-success}}

The next space station does not need to operate without government
customers to be useful. It does need an honest account of who pays,
which services are delivered and why public spending is justified.

NASA can support a commercial market by making its own needs clear and
purchasing services that meet them. Policymakers should then assess
private demand on its own evidence, rather than assuming it follows
automatically from private ownership. The test is whether the
arrangement provides dependable access to space at a defensible total
cost, while giving other customers a practical reason to participate.

\hypertarget{references}{%
\subsubsection{References}\label{references}}

\begin{enumerate}
\def\labelenumi{\arabic{enumi}.}
\tightlist
\item
  NASA. \emph{Commercial Space Frequently Asked Questions}. 7 April
  2024. Cited for NASA's stated objectives and uses of low Earth orbit.
  \href{https://www.nasa.gov/humans-in-space/leo-economy-frequently-asked-questions/}{NASA
  source}.
\item
  U.S. Government Accountability Office. \emph{Low-Earth Orbit: NASA
  Faces Impending Decisions for Replacing International Space Station
  with Commercial Stations}. GAO-26-107805, 17 June 2026.
  \href{https://www.gao.gov/products/gao-26-107805}{GAO report}.
\item
  Dooren, J. M. \emph{NASA Seeks Industry Input on Second Phase of
  Commercial Space Stations}. NASA, 6 July 2026. Cited as a dated
  announcement of a draft procurement approach.
  \href{https://www.nasa.gov/humans-in-space/commercial-space/nasa-seeks-industry-input-on-second-phase-of-commercial-space-stations/}{NASA
  announcement}.
\end{enumerate}

Sources consulted 6 October 2026. Recommendations and hypothetical
examples are the authors' analysis; no company revenue estimates are
presented.

\hypertarget{about-the-authors}{%
\subsubsection{About the authors}\label{about-the-authors}}

Adhvaith Ananth and Renish Joshi are students at the Texas Academy of
Mathematics and Science at the University of North Texas. Adhvaith's
interests include aerospace engineering and public policy. Renish
studies computer science. They are interested in how technical decisions
and public investment shape space activities.

\end{document}